\section{Udostępnione dane o szczegółach strony}
Pliki źródłowe strony nie są w żaden sposób zabezpieczone, ani zminifikowane, więc w bardzo łatwy sposób można je przeglądać. Jak się okazuje zawarte jest w nich kilka interesujących ścieżek:
\begin{itemize}
    \item Podglądanie koszyków
    
    Jedną ze znalezionych końcówek jest /basket/{id}. Okazuje się, że każdy użytkownik może podejrzeć dowolny koszyk pod warunkiem, że zapytanie będzie zawierać nagłówek z autoryzacją. Oznacza to, że wystarczy się zalogować, by podejrzeć wszystkie koszyki co jest oczywistym naruszeniem prywatności.

    \item Panel administracyjny
    
    Po zalogowaniu się na konto administratora otrzymuje się dostęp do specjalnego panelu pod ścieżką /\#/administration. Informacja o tej końcówce została znalezione w pliku strony main.js (Line 1, Column 29862). Trzeba jednak zaznaczyć, że jest to dość specyficzna podatność, bo jest to element podpowiedzi od twórców JuiceShop.

    \item Inne ciekawe znalezione końcówki, których nie udało się wykorzystać
    
    /wallet-web3, /web3-sandbox
\end{itemize}

{\rule{\linewidth}{0.2mm}}
\begin{center}
    \textcolor{red}{APLIKACJA UDOSTĘPNIA DANE NIEODPOWIEDNIM UŻYTKOWNIKOM}
\end{center} 
